\documentclass{KPS_assignment}
% ---------For Persian------------
% \newcommand*{\name}{کیوان جمالی}
% \newcommand*{\id}{403209743}
% \newcommand*{\professor}{دکتر شفاهی}
% \newcommand*{\cdate}{2024/12/26}
% \newcommand*{\course}{درس (شماره)}
% \newcommand*{\assignment}{تمرین شماره1}
% ----------For English------------
\newcommand*{\name}{Keivan Jamali}
\newcommand*{\id}{403209743}
\newcommand*{\professor}{Dr. Shafahi}
\newcommand*{\cdate}{1405/02/24}
\newcommand*{\course}{Airport (20\_582)}
\newcommand*{\assignment}{Homework 3}
% ----------START------------------
\usepackage{pgfplots}
\pgfplotsset{compat=1.18}
\begin{document}
\maketitle
\section*{Problem 1}

The characteristics of a new aircraft are given in Table B-1. In domestic route service, the applicable regulations require only a 0.75-h fuel reserve for this aircraft. The aircraft has an average route speed of 500 mi/h and an average fuel burn of 13 lb/mi.

\begin{table}[h]
\centering
\begin{tabular}{lc}
\hline
\textbf{Characteristic} & \textbf{Value (lb)} \\
\hline
Maximum ramp weight & 141000 \\
Maximum structural takeoff weight & 140000 \\
Maximum structural landing weight & 128000 \\
Zero-fuel weight & 112000 \\
Operating empty weight & 73274 \\
Fuel capacity & 39162 \\
Maximum structural payload & 38726 \\
Maximum passenger capacity & 27000 \\
Maximum cargo hold capacity & 12500 \\
\hline
\end{tabular}
\caption{Preliminary aircraft characteristics}
\end{table}

\begin{enumerate}
    \item[(a)] Plot payload versus range for this aircraft.
    \item[(b)] Determine the takeoff weight of this aircraft for a 1200-mi trip at maximum payload.
\end{enumerate}

\subsection*{Solution}

The reserve flight time is:

\[
t_{reserve} = 0.75 \text{ h}
\]

Since the aircraft speed is:

\[
V = 500 \text{ mi/h}
\]

the reserve distance becomes:

\[
R_{reserve} = 0.75 \times 500 = 375 \text{ mi}
\]

Given the fuel burn rate:

\[
13 \text{ lb/mi}
\]

the reserve fuel weight is:

\[
W_{fuel,reserve} = 375 \times 13 = 4875 \text{ lb}
\]

The takeoff weight relation is:

\[
W_{TO} = W_{OEW} + W_{payload} + W_{fuel}
\]

where:

\[
W_{OEW} = 73274 \text{ lb}
\]

and the maximum payload is:

\[
W_{payload,max} = 38726 \text{ lb}
\]

\subsection*{(a) Payload--Range Diagram}

At maximum payload:

\[
W_{fuel} = 140000 - 73274 - 38726
\]

\[
W_{fuel} = 28000 \text{ lb}
\]

The usable trip fuel is:

\[
W_{fuel,trip} = 28000 - 4875 = 23125 \text{ lb}
\]

Thus, the range is:

\[
R = \frac{23125}{13} \approx 1779 \text{ mi}
\]

For maximum fuel capacity:

\[
W_{fuel,max} = 39162 \text{ lb}
\]

The payload allowed by MTOW becomes:

\[
W_{payload} = 140000 - 73274 - 39162
\]

\[
W_{payload} = 27564 \text{ lb}
\]

The corresponding trip fuel is:

\[
39162 - 4875 = 34287 \text{ lb}
\]

Therefore:

\[
R = \frac{34287}{13} \approx 2637 \text{ mi}
\]

\begin{table}[h]
\centering
\begin{tabular}{cc}
\hline
\textbf{Range (mi)} & \textbf{Payload (lb)} \\
\hline
0 & 38726 \\
1779 & 38726 \\
2637 & 27564 \\
2637 & 0 \\
\hline
\end{tabular}
\caption{Payload--range characteristic points}
\end{table}

\begin{figure}[h]
\centering
\includegraphics[width=0.75\textwidth]{q1.png}
\caption{Payload--range diagram}
\label{fig:q1}
\end{figure}

\subsection*{(b) Takeoff Weight for a 1200-mi Trip}

For a trip distance of:

\[
R = 1200 \text{ mi}
\]

the required trip fuel is:

\[
W_{fuel,trip} = 1200 \times 13 = 15600 \text{ lb}
\]

Adding reserve fuel:

\[
W_{fuel,total} = 15600 + 4875 = 20475 \text{ lb}
\]

The takeoff weight becomes:

\[
W_{TO} = 73274 + 38726 + 20475
\]

\[
W_{TO} = 132475 \text{ lb}
\]

Therefore, the required takeoff weight is:

\[
\boxed{132475 \text{ lb}}
\]

%%%%%%

\section*{Problem 2}

An aircraft has the performance characteristics given in Table B-2 in normal route operation. In normal use this aircraft will consume 11000 lb of fuel in 1 h of flight. The normal cruising speed is 550 mi/h. Generally, alternate destinations are located within 100 mi of planned destination, and regulations require sufficient fuel to fly to the planned destination, fly to an alternate destination, and fly for 1.25 h at the alternate destination.

\subsection*{Aircraft Data}

\begin{table}[h!]
\centering
\begin{tabular}{lc}
\hline
Characteristic & Value (lb) \\
\hline
Maximum ramp weight & 191000 \\
Maximum structural takeoff weight & 190000 \\
Maximum structural landing weight & 160000 \\
Zero-fuel weight & 140000 \\
Operating empty weight & 100000 \\
Fuel capacity & 55000 \\
Maximum structural payload & 45000 \\
Maximum passenger capacity & 39000 \\
Maximum cargo hold capacity & 15000 \\
\hline
\end{tabular}
\caption{Aircraft performance characteristics}
\end{table}

The fuel consumption is:

\[
11000 \text{ lb/h}
\]

The cruise speed is:

\[
550 \text{ mi/h}
\]

Thus, the fuel burn per mile is:

\[
\frac{11000}{550}=20 \text{ lb/mi}
\]

The reserve fuel consists of:

\begin{itemize}
    \item Fuel to fly 100 mi to alternate airport:
\[
100 \times 20 = 2000 \text{ lb}
\]

    \item Fuel for 1.25 h holding:
\[
1.25 \times 11000 = 13750 \text{ lb}
\]
\end{itemize}

Therefore, total reserve fuel is:

\[
2000 + 13750 = 15750 \text{ lb}
\]

\subsection*{(a) Maximum Passenger Payload with 15000 lb Cargo}

The maximum payload allowed by the zero-fuel-weight limit is:

\[
140000 - 100000 = 40000 \text{ lb}
\]

Since 15000 lb of cargo is carried:

\[
\text{Passenger payload}
=
40000 - 15000
=
25000 \text{ lb}
\]

Assuming 200 lb per passenger:

\[
\frac{25000}{200}=125
\]

Therefore:

\[
\boxed{25000 \text{ lb passenger payload}}
\]

or

\[
\boxed{125 \text{ passengers}}
\]

\subsection*{(b) Fuel Required for a 2000-mi Trip}

Trip fuel:

\[
2000 \times 20 = 40000 \text{ lb}
\]

Required takeoff fuel:

\[
40000 + 15750 = 55750 \text{ lb}
\]

Including 1000 lb taxi fuel:

\[
55750 + 1000 = 56750 \text{ lb}
\]

Therefore:

\[
\boxed{56750 \text{ lb}}
\]

Since the fuel capacity is only 55000 lb, this flight is not feasible nonstop under the given assumptions.

\subsection*{(c) Additional Airborne Time Before Landing}

The aircraft carries 180 passengers and no cargo except baggage.

Assuming 200 lb per passenger:

\[
180 \times 200 = 36000 \text{ lb}
\]

Zero-fuel weight becomes:

\[
100000 + 36000 = 136000 \text{ lb}
\]

Maximum allowable fuel at takeoff:

\[
190000 - 136000 = 54000 \text{ lb}
\]

Fuel burned in the first 30 minutes:

\[
0.5 \times 11000 = 5500 \text{ lb}
\]

Remaining fuel:

\[
54000 - 5500 = 48500 \text{ lb}
\]

Aircraft weight after 30 minutes:

\[
136000 + 48500 = 184500 \text{ lb}
\]

The maximum landing weight is:

\[
160000 \text{ lb}
\]

Excess weight:

\[
184500 - 160000 = 24500 \text{ lb}
\]

This fuel must be burned before landing.

Additional airborne time required:

\[
\frac{24500}{11000}
=
2.23 \text{ h}
\]

Therefore:

\[
\boxed{2.23 \text{ h}}
\]

which is approximately:

\[
\boxed{2 \text{ h } 14 \text{ min}}
\]
\subsection*{(d) Payload--Range Diagram}

The payload--range diagram is constructed using the aircraft weight limitations and fuel capacity constraints.

\subsubsection*{Point A: Maximum Payload Condition}

The maximum payload is limited by the zero-fuel weight:

\[
W_{payload,\max}
=
140000-100000
=
40000 \text{ lb}
\]

Thus point A is:

\[
A=(0,\ 40000)
\]

\subsubsection*{Point B: Maximum Payload at MTOW}

At maximum payload, the takeoff weight is:

\[
W_{TO}
=
100000+40000+W_f
\]

Using the maximum takeoff weight:

\[
190000
=
100000+40000+W_f
\]

Thus:

\[
W_f
=
50000 \text{ lb}
\]

The reserve fuel requirement is:

\[
W_{res}
=
15750 \text{ lb}
\]

Therefore usable cruise fuel is:

\[
W_{cruise}
=
50000-15750
=
34250 \text{ lb}
\]

Since fuel burn is:

\[
20 \text{ lb/mi}
\]

the corresponding range is:

\[
R_B
=
\frac{34250}{20}
=
1712.5 \text{ mi}
\]

Thus point B is:

\[
B=(1712.5,\ 40000)
\]

\subsubsection*{Point C: Maximum Fuel Capacity}

At maximum fuel capacity:

\[
W_f
=
55000 \text{ lb}
\]

The allowable payload becomes:

\[
W_{payload}
=
190000-100000-55000
=
35000 \text{ lb}
\]

Usable cruise fuel is:

\[
55000-15750
=
39250 \text{ lb}
\]

Thus the range is:

\[
R_C
=
\frac{39250}{20}
=
1962.5 \text{ mi}
\]

Therefore point C is:

\[
C=(1962.5,\ 35000)
\]

\subsubsection*{Point D: Ferry Range}

Below point C, payload can continue decreasing while fuel remains fixed at maximum fuel capacity. Therefore the range remains constant in this simplified model.

Thus the ferry-range point is:

\[
D=(1962.5,\ 0)
\]

\begin{figure}[h!]
\centering
\includegraphics[width=0.85\textwidth]{q2.png}
\caption{Payload versus route range diagram}
\label{fig:q2}
\end{figure}

%%%%%%%%%%%%%%

\section*{Problem 3}

A turbine-powered aircraft has the following runway performance characteristics:

\begin{itemize}
    \item For a normal takeoff, the liftoff distance is 6500 ft, and the aircraft reaches a height of 35 ft above the end of the runway at a distance of 7000 ft from the runway threshold.
    
    \item For an engine-failure takeoff, the liftoff distance is 7200 ft, and the aircraft reaches a height of 35 ft above the end of the runway at a distance of 7900 ft from the runway threshold.
    
    \item For an aborted takeoff, the accelerate-stop distance is 9700 ft.
    
    \item For landing, the stop distance is 4600 ft.
\end{itemize}

Determine the minimum FAR runway length requirements for this aircraft, indicating:

\begin{itemize}
    \item Length of the full-strength pavement
    \item Length of the stopway
    \item Length of the clearway
\end{itemize}

Assume the runway is to be used by this aircraft in both directions.

\subsection*{Solution}

\subsubsection*{1. Normal Takeoff}

Given:

\[
LOD_1 = 6500 \ \text{ft}
\]

\[
D_{35} = 7000 \ \text{ft}
\]

The takeoff distance is:

\[
TOD_1 = 1.15 \times D_{35}
\]

\[
TOD_1 = 1.15 \times 7000 = 8050 \ \text{ft}
\]

The maximum clearway length is:

\[
CL_{1,\max} = \frac{1}{2}\left(TOD_1 - 1.15 \times LOD_1\right)
\]

\[
CL_{1,\max}
= \frac{1}{2}\left(8050 - 1.15 \times 6500\right)
\]

\[
CL_{1,\max}
= \frac{1}{2}(8050 - 7475)
\]

\[
CL_{1,\max} = 287.5 \ \text{ft}
\]

Therefore, the takeoff run required is:

\[
TOR_1 = TOD_1 - CL_{1,\max}
\]

\[
TOR_1 = 8050 - 287.5
\]

\[
TOR_1 = 7762.5 \ \text{ft}
\]

\subsubsection*{2. Engine-Failure Takeoff}

Given:

\[
LOD_2 = 7200 \ \text{ft}
\]

\[
D_{35} = 7900 \ \text{ft}
\]

For engine-failure takeoff:

\[
TOD_2 = D_{35} = 7900 \ \text{ft}
\]

The maximum clearway length is:

\[
CL_{2,\max} = \frac{1}{2}(TOD_2 - LOD_2)
\]

\[
CL_{2,\max} = \frac{1}{2}(7900 - 7200)
\]

\[
CL_{2,\max} = 350 \ \text{ft}
\]

Thus, the takeoff run required becomes:

\[
TOR_2 = TOD_2 - CL_{2,\max}
\]

\[
TOR_2 = 7900 - 350
\]

\[
TOR_2 = 7550 \ \text{ft}
\]

\subsubsection*{3. Aborted Takeoff}

Given:

\[
DAS = 9700 \ \text{ft}
\]

\subsubsection*{4. Landing Requirement}

Given stopping distance:

\[
SD = 4600 \ \text{ft}
\]

Using FAR landing requirement:

\[
LD = \frac{SD}{0.6}
\]

\[
LD = \frac{4600}{0.6}
\]

\[
LD = 7667 \ \text{ft}
\]

\subsubsection*{5. Runway Length Determination}

The required full runway length is:

\[
FL = \max(TOD_1, TOD_2, DAS, LD)
\]

\[
FL = \max(8050, 7900, 9700, 7667)
\]

\[
FL = 9700 \ \text{ft}
\]

The required full-strength pavement length is:

\[
FS = \max(TOR_1, TOR_2, LD)
\]

\[
FS = \max(7762.5, 7550, 7667)
\]

\[
FS = 7762.5 \ \text{ft}
\]

The required stopway length is:

\[
SW = DAS - \max(TOR_1, TOR_2, LD)
\]

\[
SW = 9700 - 7762.5
\]

\[
SW = 1937.5 \ \text{ft}
\]

The required clearway length is:

\[
CL = \min(FL - DAS, CL_{1,\max}, CL_{2,\max})
\]

\[
CL = \min(9700 - 9700, 287.5, 350)
\]

\[
CL = 0 \ \text{ft}
\]

\subsubsection*{Final Answers}

\[
\boxed{
\text{Full runway length} = 9700 \ \text{ft}
}
\]

\[
\boxed{
\text{Full-strength pavement length} = 7762.5 \ \text{ft}
}
\]

\[
\boxed{
\text{Stopway length} = 1937.5 \ \text{ft}
}
\]

\[
\boxed{
\text{Clearway length} = 0 \ \text{ft}
}
\]

%%%%%%%%%%%%%%

\section*{Problem 4}

An airline is considering the use of a particular aircraft on a non-stop route of 1230 mi between two airports. The nearest alternate airport at either destination is 260 mi away. The time fuel reserves are included in the item listed as ``typical operating empty weight plus time fuel reserves'' in Table~B-3. Table~B-4 indicates the takeoff requirements for this aircraft.

The airport under study has a runway of 8000 ft available with a maximum difference in centerline elevation of 15 ft. The airport elevation is 500 ft above mean sea level (AMSL), and the mean daily highest temperature in the hottest month is $90^\circ$F.

\begin{enumerate}
    \item[(a)] Find the maximum allowable takeoff weight of this aircraft at the study airport.

    \item[(b)] If the aircraft takes off from the other airport with a takeoff weight of 338{,}000 lb with the maximum allowable payload and minimum fuel, find:
    \begin{enumerate}
        \item the maximum allowable payload at the origin airport,
        \item the landing weight,
        \item and the landing runway requirements at the destination airport.
    \end{enumerate}
\end{enumerate}

\bigskip

\begin{table}[h!]
\centering
\caption{Landing Runway Requirements}
\small
\begin{tabular}{c|ccccccc}
\hline
\multicolumn{8}{c}{\textbf{Maximum Allowable Landing Weight, 1000 lb}} \\
\multicolumn{8}{c}{Airport elevation, ft} \\
Temp. ($^\circ$F) & 0 & 1000 & 2000 & 3000 & 4000 & 5000 & 6000 \\
\hline
80  & 363.5 & 363.5 & 363.5 & 363.5 & 363.5 & 363.5 & 363.5 \\
85  & 363.5 & 363.5 & 363.5 & 363.5 & 363.5 & 363.5 & 359.8 \\
90  & 363.5 & 363.5 & 363.5 & 363.5 & 363.5 & 363.5 & 350.6 \\
95  & 363.5 & 363.5 & 363.5 & 363.5 & 363.5 & 354.6 & 341.7 \\
100 & 363.5 & 363.5 & 363.5 & 363.5 & 358.4 & 345.5 & 332.8 \\
\hline
\end{tabular}

\vspace{0.5cm}

\begin{tabular}{c|ccccccc}
\hline
\multicolumn{8}{c}{\textbf{Landing Runway Length, 1000 ft}} \\
\multicolumn{8}{c}{Airport elevation, ft} \\
Weight (1000 lb) & 0 & 1000 & 2000 & 3000 & 4000 & 5000 & 6000 \\
\hline
280 & 5.44 & 5.57 & 5.71 & 5.85 & 6.00 & 6.15 & 6.31 \\
290 & 5.59 & 5.73 & 5.87 & 6.02 & 6.17 & 6.33 & 6.49 \\
300 & 5.75 & 5.89 & 6.04 & 6.19 & 6.35 & 6.51 & 6.68 \\
310 & 5.90 & 6.05 & 6.20 & 6.36 & 6.52 & 6.69 & 6.86 \\
320 & 6.06 & 6.21 & 6.37 & 6.53 & 6.70 & 6.87 & 7.05 \\
330 & 6.21 & 6.37 & 6.54 & 6.70 & 6.87 & 7.05 & 7.23 \\
340 & 6.37 & 6.53 & 6.69 & 6.86 & 7.03 & 7.21 & 7.40 \\
350 & 6.51 & 6.66 & 6.83 & 7.00 & 7.19 & 7.37 & 7.57 \\
360 & 6.64 & 6.81 & 6.98 & 7.16 & 7.35 & 7.54 & 7.75 \\
370 & 6.77 & 6.96 & 7.14 & 7.33 & 7.52 & 7.71 & 7.92 \\
\hline
\end{tabular}

\vspace{0.5cm}

\begin{flushleft}
Maximum structural takeoff weight $= 408{,}900$ lb \\
Maximum structural landing weight $= 363{,}500$ lb \\
Typical operating weight plus time fuel reserves $= 263{,}300$ lb \\
Average fuel consumption $= 30$ lb/mi \\
Typical maximum passenger load $= 54{,}000$ lb \\
Maximum structural payload $= 99{,}660$ lb \\
Fuel capacity $= 21{,}762$ gal
\end{flushleft}

\end{table}

\newpage

\begin{table}[h!]
\centering
\caption{Takeoff Runway Requirements}
\small

\begin{tabular}{c|ccccccc}
\hline
\multicolumn{8}{c}{Airport elevation, ft} \\
Temp. ($^\circ$F) & 0 & 1000 & 2000 & 3000 & 4000 & 5000 & 6000 \\
\hline
\multicolumn{8}{c}{\textbf{Maximum Allowable Takeoff Weight, 1000 lb}} \\
80  & 408.9 & 400.6 & 392.9 & 382.8 & 369.1 & 356.2 & 343.9 \\
85  & 409.0 & 400.6 & 386.0 & 372.3 & 359.4 & 347.1 & 335.4 \\
90  & 403.6 & 388.8 & 374.9 & 362.0 & 349.7 & 338.1 & 326.9 \\
95  & 391.1 & 377.3 & 364.2 & 352.0 & 340.3 & 329.2 & 318.4 \\
100 & 379.6 & 366.4 & 354.1 & 342.4 & 331.2 & 320.4 & 309.9 \\
\hline
\end{tabular}

\vspace{0.5cm}

\begin{tabular}{c|ccccccc}
\hline
\multicolumn{8}{c}{\textbf{Reference Factor}} \\
Temp. ($^\circ$F) & 0 & 1000 & 2000 & 3000 & 4000 & 5000 & 6000 \\
\hline
80  & 46.0 & 48.2 & 50.5 & 53.4 & 57.0 & 61.0 & 65.4 \\
85  & 46.3 & 48.5 & 51.8 & 55.2 & 59.0 & 63.1 & 67.8 \\
90  & 47.2 & 50.3 & 53.6 & 57.2 & 61.1 & 65.5 & 70.4 \\
95  & 48.8 & 52.0 & 55.4 & 59.2 & 63.3 & 68.0 & 73.2 \\
100 & 50.5 & 53.8 & 57.4 & 61.4 & 65.8 & 70.7 & 76.3 \\
\hline
\end{tabular}

\vspace{0.5cm}

\begin{tabular}{c|ccccc}
\hline
\multicolumn{6}{c}{\textbf{Takeoff Runway Length, 1000 ft}} \\
\multicolumn{6}{c}{Reference factor} \\
Weight (1000 lb) & 40 & 50 & 60 & 70 & 80 \\
\hline
320 & 4.23 & 5.28 & 6.39 & 7.50 & 8.56 \\
330 & 4.44 & 5.58 & 6.78 & 7.98 & 9.09 \\
340 & 4.67 & 5.89 & 7.18 & 8.46 & 9.63 \\
350 & 4.90 & 6.21 & 7.60 & 8.95 & 10.16 \\
360 & 5.13 & 6.55 & 8.03 & 9.45 & 10.68 \\
370 & 5.38 & 6.89 & 8.46 & 9.94 & -- \\
380 & 5.63 & 7.25 & 8.91 & 10.44 & -- \\
390 & 5.89 & 7.63 & 9.36 & 10.92 & -- \\
400 & 6.17 & 8.01 & 9.82 & -- & -- \\
\hline
\end{tabular}

\end{table}

\subsection{Solution}

Given data:

\begin{align*}
R &= 1230 \text{ mi} \\
R_{alt} &= 260 \text{ mi} \\
h &= 500 \text{ ft} \\
T &= 90^\circ F \\
L_{avail} &= 8000 \text{ ft}
\end{align*}

The runway centerline elevation difference is:

\[
\Delta h = 15 \text{ ft}
\]

Therefore, the runway slope is:

\[
S=\frac{15}{8000}\times100=0.1875\%
\]

Since the slope is very small, no correction is required.

\subsubsection*{(a) Maximum Allowable Takeoff Weight}

Using Table B-4 for a temperature of $90^\circ F$:

\begin{center}
\begin{tabular}{c|c}
Elevation (ft) & Maximum Takeoff Weight (1000 lb) \\
\hline
0 & 403.6 \\
1000 & 388.8
\end{tabular}
\end{center}

Interpolating for 500 ft elevation:

\[
W_{TO,max}
=
403.6
+
\frac{500}{1000}(388.8-403.6)
\]

\[
W_{TO,max}=396.2 \text{ klb}
\]

Thus,

\[
\boxed{
W_{TO,max}\approx396,200 \text{ lb}
}
\]

\subsubsection*{(b) Payload, Landing Weight, and Landing Runway Requirement}

The total distance including the alternate airport is:

\[
R_{tot}=1230+260=1490 \text{ mi}
\]

The average fuel consumption is:

\[
30 \text{ lb/mi}
\]

Therefore, the required fuel weight is:

\[
W_f=1490\times30
\]

\[
W_f=44,700 \text{ lb}
\]

The operating weight plus reserves is:

\[
W_{OE}=263,300 \text{ lb}
\]

The payload corresponding to the maximum allowable takeoff weight is:

\[
W_P
=
W_{TO,max}-W_{OE}-W_f
\]

\[
W_P
=
396,200-263,300-44,700
\]

\[
W_P=88,200 \text{ lb}
\]

However, the aircraft limitations are:

\begin{align*}
\text{Typical maximum passenger load} &= 54,000 \text{ lb} \\
\text{Maximum structural payload} &= 99,660 \text{ lb}
\end{align*}

Therefore, the allowable payload is limited to:

\[
\boxed{
W_P=54,000 \text{ lb}
}
\]

\subsubsection*{Landing Weight}

The actual takeoff weight becomes:

\[
W_{TO}
=
263,300+54,000+44,700
\]

\[
W_{TO}=362,000 \text{ lb}
\]

Fuel burned during the trip from origin to destination:

\[
W_{burn}=1230\times30
\]

\[
W_{burn}=36,900 \text{ lb}
\]

Therefore, the landing weight is:

\[
W_L
=
362,000-36,900
\]

\[
W_L=325,100 \text{ lb}
\]

Thus,

\[
\boxed{
W_L=325,100 \text{ lb}
}
\]

This value is below the maximum structural landing weight:

\[
325,100 < 363,500 \text{ lb}
\]

Therefore, the landing weight is acceptable.

\subsubsection*{Landing Runway Requirement}

Using Table B-3:

For 320 klb:

\[
L_{320}
=
6.06
+
\frac{500}{1000}(6.21-6.06)
\]

\[
L_{320}=6.135 \text{ kft}
\]

For 330 klb:

\[
L_{330}
=
6.21
+
\frac{500}{1000}(6.37-6.21)
\]

\[
L_{330}=6.29 \text{ kft}
\]

Interpolating for $325.1$ klb:

\[
L
=
6.135
+
\frac{325.1-320}{330-320}(6.29-6.135)
\]

\[
L\approx6.21 \text{ kft}
\]

Therefore,

\[
\boxed{
L_{landing}\approx6210 \text{ ft}
}
\]

\end{document}

